% sec4_4.tex — Section 4.4: Single-Equation versus Multiple-Equation Estimation

\section{Single-Equation versus Multiple-Equation Estimation}
\label{sec:4.4}

An obvious alternative to multiple-equation GMM, which estimates the stacked
coefficient vector $\boldsymbol{\delta}$ $(\sum_m L_m \times 1)$ jointly, is to apply
the single-equation GMM separately to each equation and then stack the estimated
coefficients. If the weighting matrix in the single-equation GMM for the $m$-th
equation is $\hat{\mathbf{W}}_{mm}$ $(K_m \times K_m)$, the resulting equation-by-equation
GMM estimator of the stacked coefficient vector $\boldsymbol{\delta}$ can be written as
the multiple-equation GMM estimator $\hat{\boldsymbol{\delta}}(\hat{\mathbf{W}})$ given
in~\eqref{eq:4.2.3}, where the $\sum_m K_m \times \sum_m K_m$ weighting matrix
$\hat{\mathbf{W}}$ is a block diagonal matrix whose $m$-th block is
$\hat{\mathbf{W}}_{mm}$:
\begin{equation}
\hat{\mathbf{W}} = \begin{bmatrix}
\hat{\mathbf{W}}_{11} & & \\
& \ddots & \\
& & \hat{\mathbf{W}}_{MM}
\end{bmatrix}.
\label{eq:4.4.1}
\end{equation}
This is because both $\mathbf{S}_{\mathbf{xz}}$
$(\sum_m K_m \times \sum_m L_m)$ and $\hat{\mathbf{W}}$ are block diagonal.
(The reader should actually substitute the expressions~\eqref{eq:4.2.2}
for $\mathbf{S}_{\mathbf{xz}}$ and~\eqref{eq:4.4.1} for $\hat{\mathbf{W}}$
into~\eqref{eq:4.2.3} to verify that
$\mathbf{S}_{\mathbf{xz}}' \hat{\mathbf{W}} \mathbf{S}_{\mathbf{xz}}$ and hence
$(\mathbf{S}_{\mathbf{xz}}' \hat{\mathbf{W}} \mathbf{S}_{\mathbf{xz}})^{-1}
\mathbf{S}_{\mathbf{xz}}' \hat{\mathbf{W}}$ are block diagonal,
using formulas (A.5) and (A.9) of Appendix~A. Set $M = 2$ if it makes it easier.)

Therefore, the difference between the multiple-equation GMM and the equation-by-equation
GMM estimation of the stacked coefficient vector $\boldsymbol{\delta}$ lies in the
choice of the $\sum_m K_m \times \sum_m K_m$ weighting matrix $\hat{\mathbf{W}}$.
Put differently, the equation-by-equation GMM is a particular multiple-equation GMM.

\subsection*{When Are They ``Equivalent''?}

As seen in Chapter~3, if the equation is just identified, the choice of the weighting
matrix $\hat{\mathbf{W}}$ does not matter because the GMM estimator, regardless
of the choice of $\hat{\mathbf{W}}$, numerically equals the IV estimator.
The same is true for multiple-equation GMM: if each equation of the system is exactly
identified so that $L_m = K_m$ for all $m$, then $\mathbf{S}_{\mathbf{xz}}$
in the GMM formula~\eqref{eq:4.2.3} is a square matrix and the GMM estimator
for any choice of $\hat{\mathbf{W}}$ becomes the multiple-equation IV estimator
$\mathbf{S}_{\mathbf{xz}}^{-1} \mathbf{s}_{\mathbf{xy}}$
(and, since $\mathbf{S}_{\mathbf{xz}}$ is block diagonal, the estimator is just a
collection of single-equation IV estimators).

On the other hand, if at least one of the $M$ equations is overidentified, the choice
of $\hat{\mathbf{W}}$ $(\sum_m K_m \times \sum_m K_m)$ does affect the numerical
value of the GMM estimator. Consider the efficient equation-by-equation GMM
estimator of $\boldsymbol{\delta}$ $(\sum_m L_m \times 1)$, where each equation is
estimated by the efficient single-equation GMM with an optimal choice of
$\hat{\mathbf{W}}_{mm}$ $(m = 1, 2, \ldots, M)$.
Is it as efficient as the efficient multiple-equation GMM?
The equation-by-equation estimator can be written as
$\hat{\boldsymbol{\delta}}(\hat{\mathbf{W}})$, where the
$\sum_m K_m \times \sum_m K_m$ block diagonal weighting matrix $\hat{\mathbf{W}}$
satisfies
\begin{equation}
\plim_{n \to \infty} \hat{\mathbf{W}} = \begin{bmatrix}
\bigl(\E(\varepsilon_{i1}^2\, \mathbf{x}_{i1} \mathbf{x}_{i1}')\bigr)^{-1}
& & \\
& \ddots & \\
& & \bigl(\E(\varepsilon_{iM}^2\, \mathbf{x}_{iM} \mathbf{x}_{iM}')\bigr)^{-1}
\end{bmatrix}.
\label{eq:4.4.2}
\end{equation}
Because this plim is not necessarily equal to $\mathbf{S}^{-1}$
(unless~\eqref{eq:4.4.3} below holds), the estimator is generally less efficient than
the efficient multiple-equation GMM estimator
$\hat{\boldsymbol{\delta}}(\widehat{\mathbf{S}}^{-1})$, where
$\widehat{\mathbf{S}}$ is a consistent estimator of $\mathbf{S}$.
But if the equations are ``unrelated''\footnote{Under conditional homoskedasticity,
this condition becomes the more familiar condition that
$\E(\varepsilon_{im}\, \varepsilon_{ih}) = 0$.
See the next section.} in the sense that
\begin{equation}
\E(\varepsilon_{im}\, \varepsilon_{ih}\, \mathbf{x}_{im} \mathbf{x}_{ih}') = \mathbf{0}
\quad \text{for all } m \neq h \;(= 1, 2, \ldots, M),
\label{eq:4.4.3}
\end{equation}
then $\mathbf{S}$ becomes block diagonal and for the weighting matrix
$\hat{\mathbf{W}}$ defined above
$\plim \hat{\mathbf{W}} = \mathbf{S}^{-1}$.
Since in this case
$\hat{\mathbf{W}} - \widehat{\mathbf{S}}^{-1} \pto \mathbf{0}$, we have
(as we proved in Review Question~3 of Section~3.5)
\[
\sqrt{n}\, \hat{\boldsymbol{\delta}}(\hat{\mathbf{W}})
- \sqrt{n}\, \hat{\boldsymbol{\delta}}(\widehat{\mathbf{S}}^{-1}) \pto \mathbf{0}.
\]

It then follows from Lemma~2.4(a) that the asymptotic distribution of
$\hat{\boldsymbol{\delta}}(\hat{\mathbf{W}})$ is the same as that of
$\hat{\boldsymbol{\delta}}(\widehat{\mathbf{S}}^{-1})$, so the efficient
equation-by-equation GMM estimator is an efficient multiple-equation GMM estimator.

We summarize the discussion so far as

\begin{proposition}[equivalence between single-equation and multiple-equation GMM]
\label{prop:4.3}
\leavevmode
\begin{enumerate}[label=(\alph*)]
\item \textit{If all equations are just identified, then the equation-by-equation GMM
and the multiple-equation GMM are numerically the same and equal to the IV estimator.}

\item \textit{If at least one equation is overidentified but the equations are
``unrelated'' in the sense of\/~\eqref{eq:4.4.3}, then the efficient
equation-by-equation GMM and the efficient multiple-equation GMM are
asymptotically equivalent in that $\sqrt{n}$ times the difference converges to
zero in probability.}
\end{enumerate}
\end{proposition}

For those two cases, there is no efficiency gain from joint estimation. Furthermore,
if it is known \textit{a priori} that the equations are ``unrelated,'' then the
equation-by-equation estimation should be preferred, because exploiting the
\textit{a priori} knowledge by requiring the off-diagonal blocks of $\hat{\mathbf{W}}$
to be zeros (which is what the equation-by-equation estimation entails) will most
likely improve the finite-sample properties of the estimator.

\subsection*{Joint Estimation Can Be Hazardous}

Except for cases (a) and (b), joint estimation is asymptotically more efficient. Even
if you are interested in estimating one particular equation (say, the $LW$ equation
in Example~4.1), you can generally gain asymptotic efficiency by combining it
with some other equations (see, however, Analytical Exercise~10 for an example
where joint estimation entails no efficiency gain even if the added equations are not
unrelated).

There are caveats, however. First of all, the small-sample properties of the
coefficient estimates of the equation in question might be better without joint estimation.
Second, the asymptotic result presumes that the model is \textbf{correctly specified},
that is, all the model assumptions are satisfied. If the model is misspecified,
neither the single-equation GMM nor the multiple-equation GMM is guaranteed
even to be consistent. And chances of misspecification increase as you add equations
to the system.
